\section{The Database: Locks, Transactions and Migrations}

\subsection{Why PostgreSQL does the hard part}

Sixty PCs can ask for work in the same second. If two of them read ``job 7 is
free'' at the same time, both may take it. Instead of inventing a locking
scheme, the scheduler lets PostgreSQL do it with \textbf{transactions} and
\textbf{row locks}. Each request runs in one transaction: either all of its
changes are saved together, or none are.

\subsection{SELECT \dots{} FOR UPDATE SKIP LOCKED}

\texttt{SELECT ... FOR UPDATE} reads rows \emph{and locks them} until the
transaction ends. Another transaction trying to lock the same row waits.
Adding \texttt{SKIP LOCKED} changes ``wait'' into ``skip it and look at the
next row''. That one clause turns a table into a safe work queue:

\begin{lstlisting}[language=SQL]
-- Two workers run this at the same moment.
BEGIN;
SELECT id FROM jobs
 WHERE desired_state = 'run' AND accepted_attempt_id IS NULL
 ORDER BY priority DESC, created_at
 LIMIT 1
 FOR UPDATE SKIP LOCKED;     -- worker A locks job 7, worker B skips to job 8
-- ... create the attempt ...
COMMIT;                      -- lock released
\end{lstlisting}

\subsection{A real bug: locking the whole queue}

The first version of \texttt{claim()} locked \emph{every} waiting job, then
picked one in Python:

\begin{lstlisting}[language=Python]
# before: locks all queued jobs for the whole claim
stmt = (select(Job)
        .where(Job.desired_state == "run",
               Job.accepted_attempt_id.is_(None))
        .with_for_update(skip_locked=True))
jobs = (await session.scalars(stmt)).all()
\end{lstlisting}

With two PCs asking at the same moment, PC A locked all twenty queued jobs
while it decided. PC B skipped every locked row, saw an empty queue, and went
back to sleep, even though nineteen jobs were waiting. Nothing broke, so no
test noticed; the cluster just ran slower than it should.

The fix reads candidates \emph{without} locks, decides in Python which job
fits this node, and then locks \emph{only that one row}. If another node got it
first, it moves to the next candidate:

\begin{lstlisting}[language=Python]
# after (control/app/scheduler.py, simplified)
jobs = (await session.scalars(candidates_stmt.limit(200))).all()
for job in jobs:
    if not node_fits_job(node, job, envs.get(job.environment_id)):
        continue
    locked = await session.scalar(
        select(Job).where(Job.id == job.id,
                          Job.accepted_attempt_id.is_(None))
                   .with_for_update(skip_locked=True))
    if locked is None:          # another node won this one
        continue
    chosen = locked
    break
\end{lstlisting}

The test that proves it opens two database sessions, starts a claim in the
first \emph{without committing} (so its lock stays held), and checks the second
still gets the other job. On the old code the second claim returned nothing.

\subsection{Partial unique indexes}

``A node runs at most one job at a time'' is enforced by the database itself,
not only by Python:

\begin{lstlisting}[language=SQL]
CREATE UNIQUE INDEX uq_attempts_one_active_per_node
    ON job_attempts (node_id)
 WHERE observed_state IN ('leased','starting','running',
                          'uploading','cancelling','unknown');
\end{lstlisting}

A \emph{partial} index only covers rows matching its \texttt{WHERE}. Finished
attempts are not in it, so a node can have hundreds of past attempts but only
one active. If a bug ever tried to insert a second active attempt, Postgres
would reject the insert. The idempotency key uses the same trick with
\texttt{WHERE idempotency\_key IS NOT NULL}.

\paragraph{Other use cases.} ``One active subscription per user'', ``one
default address per customer'', ``one open shift per employee''.

\subsection{Talking to Postgres from Python: SQLAlchemy async}

The control plane uses SQLAlchemy 2 with the \texttt{asyncpg} driver. Tables
are Python classes:

\begin{lstlisting}[language=Python]
class JobAttempt(Base):
    __tablename__ = "job_attempts"
    id: Mapped[uuid.UUID] = mapped_column(Uuid, primary_key=True)
    job_id: Mapped[uuid.UUID] = mapped_column(ForeignKey("jobs.id"))
    fencing_token: Mapped[uuid.UUID] = mapped_column(Uuid, unique=True)
    lease_expires_at: Mapped[datetime] = mapped_column(DateTime(timezone=True))
    observed_state: Mapped[str] = mapped_column(String(32))
\end{lstlisting}

\textbf{Why async?} A long-poll claim spends most of its 25 seconds asleep.
With \texttt{async}/\texttt{await}, one process serves hundreds of waiting
agents, because a sleeping request costs almost nothing. With one thread per
request, sixty agents would need sixty threads doing nothing.

\paragraph{A trap we hit.} An async engine keeps a pool of connections that
belong to the event loop that created them. The operations tool first called
\texttt{asyncio.run()} twice in a row; the second call started a new loop and
reused pooled connections from the first, failing with ``attached to a
different loop''. The fix was to run both steps inside one
\texttt{asyncio.run()}.

\subsection{Schema migrations with Alembic}

Early on the app created tables at startup with \texttt{create\_all()} and
patched new columns in with hand-written \texttt{ALTER TABLE ... IF NOT
EXISTS}. That works until the first real deployment, then becomes dangerous:
there is no record of which changes a database already has.

\textbf{Alembic} keeps an ordered list of migration scripts and a one-row table,
\texttt{alembic\_version}, recording which one a database has reached:

\begin{verbatim}
control/migrations/versions/
  0001_baseline_schema.py      every table as it was
  0002_job_idempotency_key.py  adds jobs.idempotency_key + unique index
  0003_node_update_ring.py     adds nodes.update_ring
\end{verbatim}

\begin{lstlisting}[language=Python]
# 0003_node_update_ring.py
revision = "0003"
down_revision = "0002"

def upgrade() -> None:
    op.add_column("nodes", sa.Column("update_ring", sa.Integer(),
                  server_default=sa.text("2"), nullable=False))

def downgrade() -> None:
    op.drop_column("nodes", "update_ring")
\end{lstlisting}

Adding a column is now: edit \texttt{models.py}, run
\verb|alembic revision --autogenerate -m "..."|, review the generated
script, commit. The API applies pending migrations on startup.

\paragraph{Adopting migrations on an existing database.} Databases created the
old way have the tables but no \texttt{alembic\_version} row. On startup,
\texttt{init\_db} detects that case, applies the old column patches, and
\emph{stamps} the database at revision \texttt{0001} (records it as already
there) before upgrading. This was tested on a copy of the real development
database, then \texttt{alembic check} confirmed the schema matched the models
exactly.

\paragraph{A detail autogenerate got wrong.} The \texttt{jobs} and
\texttt{job\_attempts} tables reference each other. Autogenerate put both
foreign keys inline in \texttt{CREATE TABLE jobs}, but \texttt{job\_attempts}
did not exist yet. The baseline was edited to add those two keys at the end
with \texttt{op.create\_foreign\_key}. Always read generated migrations.

\paragraph{Other use cases.} Django migrations, Rails Active Record migrations,
Flyway and Liquibase in Java, Prisma Migrate in Node.
